\documentclass{beamer}

\mode<presentation> {%
	\usetheme[]{Boadilla}
	\usecolortheme{orchid}
	\setbeamercovered{transparent}
	\setbeamersize{text margin left=0.7cm}
  \setbeamersize{text margin right=0.7cm}
  \setbeamertemplate{enumerate item}[square]
  \setbeamertemplate{enumerate subitem}[circle]
  \setbeamertemplate{enumerate subsubitem}[>]
  \setbeamertemplate{itemize item}[square]
  \setbeamertemplate{itemize subitem}[circle]
  \setbeamertemplate{itemize subsubitem}[>]
}

% \usepackage{bm}
% \usepackage[utf8]{inputenc}
% \usepackage[spanish]{babel}
% \usepackage{times}
% \usepackage[T1]{fontenc}
% \usepackage{graphicx}
% \usepackage{multirow}
% \usepackage{pifont}
% \usepackage{rotating}
% \usepackage{colortbl}
% \usepackage[framed]{mcode}
\usepackage[utf8]{inputenc}
\usepackage{amsmath}
\usepackage{graphicx}
\usepackage{booktabs}

\definecolor{olivegreen}{rgb}{0.0,0.5,0.0}
% \definecolor{general}{rgb}{0.9 0.9 1}%
% \newrgbcolor{colorsistema}{0 1 1}
% \newrgbcolor{colorviolet}{0.9 0.8 1}

\title[Procesamiento de Señales]{Transformada de Fourier Discreta (DFT)}
\author[J. Chacón]{%[Author, Another] % (optional, use only with lots of authors)
\textit{Ingeniería Electrónica de Comunicaciones }\\
\vspace{0.6cm}
\small Jesús Chacón Sombría}

\institute[DACYA] % (optional, but mostly needed)
{
  %\inst{1}%

\footnotesize  Departamento de Arquitectura de Computadores y Automática.\\
  Universidad Complutense de Madrid
  }

\date[Transformada Fourier Discreta] % (optional)
{\footnotesize Curso 2025-2026}%{\today / Master}
\subject{Talks}

\pgfdeclareimage[height=0.75cm]{logo}{ucm.png}

\usepackage[absolute,overlay]{textpos}
\setlength{\TPHorizModule}{1mm}
\setlength{\TPVertModule}{1mm}
\newcommand{\MyLogo}{%
\begin{textblock}{14}(116.0,0.6)
  \pgfuseimage{logo}
\end{textblock}
}

\makeatletter
\numberwithin{section}{part}
\AtBeginPart{
\beamer@tocsectionnumber=0\relax
\frame{\partpage}
}
\makeatother

\begin{document}

\begin{frame}
  \titlepage
\end{frame}

\newcommand*\oldmacro{}%
\let\oldmacro\insertshorttitle%
\renewcommand*\insertshorttitle{%
 \MyLogo
  \oldmacro\hfill%
}

\AtBeginSection[]{}

\AtBeginSubsection[]{}      

\begin{frame}
  \frametitle{Periodograma}
  
  \footnotesize
  \begin{itemize}
    \item \alert{Podemos usar la DFT para estimar el espectro de una señal/sistema.}
    \[
    X(\frac{2\pi k}{N})=X[k]=\substack{N-1\\\sum\\n=0}x[n]\text{e}^{-j\frac{2\pi}{N}kn}
    \]
    \begin{itemize}
      \item La muestra $k$ corresponde a la frecuencia discreta $\Omega=\frac{2\pi}{N}k$.
      \item Si usamos un periodo de muestreo $T_s$, entonces:
    \[
    \omega=\frac{\Omega}{T_s}=\frac{2\pi k}{N T_s},f=\frac{k}{NT_s}
    \]
    \end{itemize}

    \item \alert{La densidad espectral es $\Phi_{xx}(\omega_k)\approx\frac{1}{N\cdot f_s}|X[k]|^2$.}
    \begin{itemize}
      \footnotesize
      \item En escala logarítmica (dB/Hz):
      \[
        \Phi_{xx}|_{dB} = 20\:\text{log}_{10}\left(\frac{|X[k]|}{\sqrt{N\cdot f_s}}\right)
      \]
    \end{itemize} 
    \item \alert{La potencia es: $\frac{|X[k]|^2}{N\cdot f_s}\Delta f=\frac{|X[k]|^2}{N\cdot f_s}\cdot\frac{f_s}{N}=|\frac{X[k]}{N}|^2$.}
    \begin{itemize}
      \footnotesize
      \item En escala logarítmica (dB):
      \[
        \Phi_{xx}|_{dB} = 20\:\text{log}_{10}\left(\frac{|X[k]|}{N}\right)
      \]
    \end{itemize} 
    
    

  \end{itemize}
\end{frame}



\begin{frame}
  \frametitle{Convolución circular}
  
  \footnotesize
  \begin{itemize}
    \item Tenemos dos secuencias finitas de longitud $L$ y $M$, elegimos $N=L+M-1$:
    \begin{center}
      $\color{blue}x_1[n]\color{black} \Leftrightarrow X_1[k]=\substack{N-1\\\sum\\n=0} x_1[n]e^{-j \frac{2\pi}{N}kn}
      \quad 
      \color{blue}x_2[n]\color{black} \Leftrightarrow X_2[k]=\substack{N-1\\\sum\\n=0}x_2[n]e^{-j \frac{2\pi}{N}kn}$
    \end{center}
    
    \item Definimos $\color{blue}X_3[k]=X_1[k]X_2[k]\color{black}$ y calculamos su IDFT:\scriptsize
    \[
    \begin{aligned}
      x_3[n] &= \frac{1}{N}\sum_{k=0}^{N-1} X_3[k]e^{j \frac{2\pi}{N}kn}=\frac{1}{N}\sum_{k=0}^{N-1} \left[\sum_{m=0}^{N-1}x_1[m]e^{-j \frac{2\pi}{N}km} \right] \left[\sum_{l=0}^{N-1}x_2[l]e^{-j \frac{2\pi}{N}kl} \right]e^{j \frac{2\pi}{N}kn}=\\
             &= \frac{1}{N}\sum_{m=0}^{N-1} x_1[m]  \sum_{l=0}^{N-1}x_2[l] \left[ \sum_{k=0}^{N-1}e^{j \frac{2\pi}{N}k(n-m-l)} \right ]=\\
             &= \sum_{m=0}^{N-1} x_1[m] x_2[n-m-pN]=\sum_{n=0}^{N-1} x_1[n] x_2[(m-n)|_N]=(x*y)_N[n]
    \end{aligned}
    \]

    \item \alert{$(x*y)_N[m]=\sum_{n=0}^{N-1} x_1[n] x_2[(m-n)|_N]  \:\:\:\:   \leftrightarrow \:\:\:\: X[k]Y[k]$}
    
    \rule{11cm}{0.4pt}
    
    $\substack{N-1\\\sum\\k=0}e^{j \frac{2\pi}{N}k(n-m-l)} = \left \{\begin{array}{l} N,   \:\:\: n-m-l=pN, p\in \mathcal{Z} \\ 0,  \:\:\: c.c. \end{array}\right .$  
  \end{itemize}
\end{frame}
  
  
  \begin{frame}
  \frametitle{Sistemas LTI y filtros discretos}
  
  \begin{itemize}
    \item Respuesta $y[n]$ de un sistema LTI discreto a una entrada $x[n]$:
    \begin{itemize}
      \vspace{0.2cm}
      \item \textbf{Metodo 1}:
      \begin{itemize}
        \item $y[n]=\substack{\infty\\\sum\\m=-\infty} x[m]h[n-m]$
      \end{itemize}  
      \vspace{0.2cm}
      \item \textbf{Metodo 2}:
      \begin{itemize}
        \item Transformada de Fourier: $Y(e^{jw})=H(e^{jw})X(e^{jw})$\\
        \vspace{0.1cm}
        $X(e^{jw})=\sum_{n=-\infty}^{\infty} x[n] e^{-jwn}$\\
        $H(e^{jw})=\sum_{n=-\infty}^{\infty} h[n] e^{-jwn}$
      \end{itemize}  
    \end{itemize}  
  \end{itemize}

\end{frame}

\begin{frame}
    \frametitle{Sistemas LTI y filtros discretos}
    
    \footnotesize
    \alert{¿Podemos usar la DFT para calcular la respuesta de sistemas LTI?}
  \begin{itemize}
    \item Las operaciones correspondientes en el caso de la DFT son:
    \begin{itemize}
    \footnotesize
    \item $y[n]=\substack{N-1\\\sum\\m=0} x[m]h[n-m]_N$ para señales de longitud $\le N$
    \vspace{0.1cm}
    
    \item $Y[k]=H[k]X[k]$
    \vspace{0.1cm}
    \item $X[k]=\substack{N-1\\\sum\\n=0} x[n]e^{j\frac{2\pi}{N}kn}$ para señales de longitud $\le N$
    \end{itemize}
    \item Para poder realizar las operaciones con la DFT necesitamos que:  
    \begin{itemize}
      \footnotesize
      \item $x[n]$ sea una señal de longitud finita $L< N$
      \item $h[n]$ sea una señal de longitud finita $M < N$  $\rightarrow $ sistema FIR.
      \item Para que no haya elementos del desplazamiento circular, $N \ge L+M-1$
    \end{itemize}
  \end{itemize}
\end{frame}
  
\begin{frame}
  \frametitle{Sistemas LTI y filtros discretos}
  
  \begin{itemize}
    \item \textbf{Si el sistema $h[n]$ es IIR}:
    \begin{itemize}
      \item Podemos obtener la DFT como $H[k]=H[e^{\frac{2\pi k}{N}}]$. 
      \item No podemos usar la DFT para calcular la respuesta del sistema.
    \end{itemize}
    \vspace{0.2cm}
  
    \item \textbf{Si la entrada $x[n]$ no tiene longitud finita:}
    \begin{itemize}
      \item Podemos obtener la DFT como $X[k]=X(e^{\frac{2\pi k}{N}})$.
      \item No utilizar la DFT para calcular la respuesta de un sistema LTI.
    \end{itemize}
    \vspace{0.2cm}

    \item Para $x[n],h[n]$ periódicas, caso particular:
    \begin{itemize}
      \item Señales de longitud finita $N,M$ (un periodo). 
    \end{itemize}
  \end{itemize}
\end{frame}
  
  
\begin{frame}
  \frametitle{Ejemplo de respuesta de sistema LTI}
  
  \footnotesize
  \begin{itemize}
    \item \color{olivegreen}Calcular la respuesta del sistema $h[n]=\{0,2\}$ a la entrada  $x[n]=\{1,3,4\}$.\color{gray}
    \begin{itemize}
      \item \color{black}A traves de la convolución:\color{gray}
    \end{itemize}

    \scriptsize
    $\left . \!\!\!\!\! \begin{array}{l} y[0]=\sum_{n=0}^{1}h[n]x[0-n]=0\cdot 1 + 2 \cdot 0=0\\
    y[1]=\sum_{n=0}^{1}h[n]x[1-n]=0\cdot 3 + 2 \cdot 1=2\\
    y[2]=\sum_{n=0}^{1}h[n]x[2-n]=0\cdot 4+ 2 \cdot 3=6\\
    y[3]=\sum_{n=0}^{1}h[n]x[3-n]=0\cdot 0 + 2 \cdot 4=8\end{array}\!\!\right\}$ $\rightarrow y[n]=(0,2,6,8)^T$
  
    \begin{itemize}
      \item \color{black}A través de la DFT:\color{gray}
    \end{itemize}
      
      \scriptsize 
      $N=3+2-1=4$\\[1ex]
      $x[n]=(1,3,4,0)^T ; h[n]=(0,2,0,0)^T$\\[1ex]
      $W=\begin{pmatrix} 
        1 &  1 &  1 &  1 \\
        1 & -j & -1 &  j \\
        1 & -1 &  1 & -1 \\
        1 &  j & -1 & -j \\
      \end{pmatrix}
      \quad
      W^{-1}=\begin{pmatrix} 
        1 &  1 &  1 &  1 \\
        1 &  j & -1 & -j \\
        1 & -1 &  1 & -1 \\
        1 & -j & -1 &  j \\
      \end{pmatrix}$\\[1ex]      
      $X[k]=W\cdot x[n]=\left(8, -3-3j, 2, -3+3j\right)^T$\\[1ex]
      $H[k]=W\cdot h[n]=\left(2, -2j, -2, 2j\right)^T$\\[1ex]
      $y[n]=\frac{1}{N}W^{-1}\cdot H[K]\cdot X[k]=(0,2,6,8)^T$
  \end{itemize}
\end{frame}

\begin{frame}
  \frametitle{Conceptos importantes}

  \begin{itemize}\scriptsize
    \item \alert{DFT para estimar el espectro.}
    \[
      X[k]=X(\frac{2\pi k}{N})=\substack{N-1\\\sum\\n=0}x[n]\text{e}^{-j\frac{2\pi}{N}kn}
    \]

    \item \alert{La potencia/densidad espectral es:}
    \[
      P_{x}(\omega_k)=\left|\frac{X[k]}{N}\right|^2 \quad P_{x}(\omega_k)|_{dB} = 20\:\text{log}_{10}\left(\frac{|X[k]|}{N}\right)
    \]
    \[
    \Phi_{xx}(\omega_k)=\frac{|X[k]|^2}{N\cdot f_s} \quad \Phi_{xx}(\omega_k)|_{dB} = 20\:\text{log}_{10}\left(\frac{|X[k]|}{\sqrt{N\cdot f_s}}\right)
    \]

    \item \alert{Convolución circular:}
    \[
      \substack{N-1\\\sum\\k=0}e^{j \frac{2\pi}{N}k(n-m-l)} = \left \{\begin{array}{l} N,   \:\:\: n-m-l=pN, p\in \mathcal{Z} \\ 0,  \:\:\: c.c. \end{array}\right .
    \]

    \item \alert{Respuesta LTI:}
    \begin{itemize}\scriptsize
      \item \textbf{Si el sistema $h[n]$ es IIR}:
      \begin{itemize}\scriptsize
        \item Podemos obtener la DFT con $H[k]=H[e^{\frac{2\pi k}{N}}]$. 
        \item No podemos usar la DFT para calcular la respuesta del sistema.
      \end{itemize}
      \vspace{0.2cm}
    
      \item \textbf{Si la entrada $x[n]$ no tiene longitud finita:}
      \begin{itemize}\scriptsize
        \item Podemos obtener la DFT como $X[k]=X(e^{\frac{2\pi k}{N}})$.
        \item No utilizar la DFT para calcular la respuesta de un sistema LTI.
      \end{itemize}
      \vspace{0.2cm}
    \end{itemize}
  \end{itemize}

\end{frame}

\end{document}